\newcommand{\store}[1]{\langle #1 \rangle}

\begin{frame}[fragile]
    \frametitle{Kripke's Semantics}

\textbf{Idea:} Encode the set of values $x\mapsto \{1,2,3\}$ as a set of worlds: $\{ [x\mapsto 1], [x\mapsto 2], [x\mapsto 3] \}$.

\textbf{Relation over worlds:} For underapproximation, we care about a set of worlds instead of a single world, which can be expressed as an equivalence class.

\textbf{Existing approaches:} 
\begin{itemize}
    \item Classical modal logic S5, where the relation over worlds is an equivalence relation.
    \item Epistemic logic, which includes knowledge and belief.
    \begin{itemize}
        \item The knowledge operator $K$ is a modal operator, interpreted as ``I know that'', which is modeled by an equivalence relation.
        \item The belief operator $B$ is a modal operator, interpreted as ``I believe that'', which is symmetric and transitive, but not necessarily reflexive (allowing for misunderstanding).
        \begin{align*}
            B\phi \not\models \phi
        \end{align*}
    \end{itemize}
\end{itemize}
        
\end{frame}

\begin{frame}[fragile]
    \frametitle{Kripke's Semantics}

\textbf{Idea:} Introduce a new underapproximation modality $\boxtimes$; consider the following model:

% https://tikzcd.yichuanshen.de/#N4Igdg9gJgpgziAXAbVABwnAlgFyxMJZABgBpiBdUkANwEMAbAVxiRAHcB9ARkQAJkADz4AdEQFs6aODgh9uFEAF9S6TLnyEU3Ugur1mrRB04BmFILGTpsvqcUq12PASIAmclVqMWbLm-4hKykZCDcHLxgoAHN4IlAAMwAnCHEkMhBZJB0QBiwwIxAoCCYAIwZWagALGDooNkgCkGocOiwGBoJWRxBk1KQPTIhs6tr640bK3PzC4rKK5sy2jomu5VVelLTEDKzEQbym4znyqZq6zqaW5cvuiiUgA
\begin{tikzcd}
    {w_1: [x \mapsto 1]} \arrow[rd, no head] \arrow[rr, no head] &                    & {w_2: [x\mapsto2]} \arrow[ld, no head] \\
                                                                                         & {w_3:[x\mapsto 3]} &                                                   
    \end{tikzcd}

\textbf{Example:} 
\begin{itemize}
    \item $w_1 \models 1 \leq x \leq 10$ means the value of $x$ in world $w_1$ is in the range $1$ to $10$.
    \item $w_1 \models \boxtimes (1 \leq x \leq 3)$ means all worlds that satisfy $1 \leq x \leq 3$ are equivalent to $w_1$.
    \item $w_1 \models \square (1 \leq x \leq 3)$ (the standard S5 necessity modality) means all worlds equivalent to $w_1$ satisfy $1 \leq x \leq 3$.
    \item Note that $w_1 \models \square (1 \leq x \leq 10)$ holds; however, $w_1 \models \boxtimes (1 \leq x \leq 10)$ does not hold.
\end{itemize}
 
\end{frame}

\begin{frame}[fragile]
    \frametitle{Local Reasoning}

\textbf{Idea:} We still want to enable local reasoning (frame rules and other substructural rules); that is, consider ``pieces of worlds (heaps)''.

\textbf{Problem:} The current definition of $\boxtimes$ is too strong. For example, consider $w_1 \models \boxtimes (1 \leq x \leq 3)$. 

\vspace{0.5cm}
% https://tikzcd.yichuanshen.de/#N4Igdg9gJgpgziAXAbVABwnAlgFyxMJZABgBpiBdUkANwEMAbAVxiRAHcB9ARkQAJkADz4AdEQFs6aODgh9uFEAF9S6TLnyEU3Ugur1mrRB04BmFILGTpsvqcUq12PASIAmclVqMWbLm-4hKykZCDcHVRAMZ013XS8DX2MuABYLYJs5UwBuPgBPUQkQ2wBWBy8YKABzeCJQADMAJwhxJDIQWSQdEAYsMCMQOAheqBBqAAsYOlHjSH6xjrosBjY51kcQJpakDw6ILompmfACVmpe+eMhkYWcJZXZ0+VIrdbEds7EXcnp1aeN14HPZIUyHX6IMBMBgMfQ+AYAfmUFCUQA
\begin{tikzcd}
    {w_1: [x \mapsto 1]} \arrow[rd, no head] \arrow[rr, no head] &                                   & {w_2: [x\mapsto2]} \arrow[ld, no head] \\
                                                                 & {w_3:[x\mapsto 3]} \arrow[r, "?"] & {w_4:[x\mapsto 3; y \mapsto 5]}       
    \end{tikzcd}

\vspace{0.5cm}
\textbf{Solution:} Not ``all worlds that satisfy $\phi$'', but ``all \textbf{minimal} worlds that satisfy $\phi$'', where $[x\mapsto 3; y \mapsto 5]$ is not a minimal world for $\boxtimes (1 \leq x \leq 3)$. 

\end{frame}

\begin{frame}[fragile]
    \frametitle{Better definition}

\textbf{Base (Valuation Function):} A world can be represented as a substitution; $\Dom{\sigma}$ indicates the variables in the substitution.

\textbf{Inclusion (Projection) relation:} $w_1 \subseteq w_2$ means all variables in $w_1$ have the same assignment in $w_2$; we also say $w_2$ can be projected to $w_1$.

\textbf{Equivalence relation:} There is an equivalence relation $\sim$ over worlds $w_1$ and $w_2$. We require that equivalent worlds have the same domain, that is, $\Dom{w_1} = \Dom{w_2}$. The intuition is that only aligned assignments can be grouped together, or we can distinguish the set of variables but don't care about the actual assignment.

\textbf{Preservation of projection relation:} Note that the equivalence relation preserves the projection relation: for all $w_1 \subseteq w_1'$, we have $w_2 \subseteq w_2'$.
\begin{align*}
    w_1' \sim w_2' \implies w_1 \sim w_2
\end{align*}
\noindent Note that the reverse is not true.


\end{frame}

\begin{frame}[fragile]
    \frametitle{Better definition}

Consider a model $M_1$.

% https://tikzcd.yichuanshen.de/#N4Igdg9gJgpgziAXAbVABwnAlgFyxMJZABgBpiBdUkANwEMAbAVxiRAHcB9ARkQAJkADz4AdEQFs6aODgh9uFEAF9S6TLnyEU3clVqMWbLgCZ+QsZOmzjilWux4CRMjer1mrRB04BmMwE9RCSkZOR9bVRAMB00iHVd9DyNOABYAoMtQvhSI+w0nbVIFNwNPbwA2M2ELENk+YwBuPn8aqzkc5Ujo-K1kYyK9d0MvLgB2KozauW4mluC27Nyo9UdesmLE4e8AVgnWrJnmyYXw5T0YKABzeCJQADMAJwhxJDIQWSQdEAALGDooNiQMCsOwgR7PJD9d4QJA+ai-f6AgggyLgl6IFLUD6IL4MLDAthQOhwX4A0FopCY6Gw6h4gleIkki6de5PdHbLEwxBvOllRmkllgtlIDnUxBwkC8wnEgXk4WIcqc160-F8mXMuUQhVKxBQqUM9Vk1HyxViqkIgFeIEgihKIA
\begin{tikzcd}
    {w_1: [x \mapsto 1]} \arrow[r, no head]                                                  & {w_2: [x\mapsto2]}                                                   &                                                                         \\
    {w_5: [x \mapsto 1; y \mapsto 3]} \arrow[u, dashed] \arrow[d, dashed] \arrow[r, no head] & {w_6: [x \mapsto 2; y\mapsto 4]} \arrow[u, dashed] \arrow[d, dashed] & {w_7: [x \mapsto 1; y\mapsto 4]} \arrow[llu, dashed] \arrow[ld, dashed] \\
    {w_3: [y \mapsto 3]} \arrow[r, no head]                                                  & {w_4: [y \mapsto 4]}                                                 &                                                                        
    \end{tikzcd}
    
The equivalence class $\{w_5, w_6\}$ represents execution $(\lambda x.x+2)\;\{1,2\}$.
\begin{itemize}
    \item After projection, $\{w_1, w_2\}$ is still a valid equivalence class, indicating the input $\{1,2\}$.
    \item After projection, $\{w_3, w_4\}$ is still a valid equivalence class, indicating the output $\{3,4\}$.
    \item The set $\{w_5, w_7\}$ is not a valid equivalence class, although $w_3 \subseteq w_5$ and $w_4 \subseteq w_7$.
\end{itemize}

\end{frame}

\begin{frame}[fragile]
    \frametitle{Better definition}

% https://tikzcd.yichuanshen.de/#N4Igdg9gJgpgziAXAbVABwnAlgFyxMJZABgBpiBdUkANwEMAbAVxiRAHcB9ARkQAJkADz4AdEQFs6aODgh9uFEAF9S6TLnyEU3clVqMWbLgCZ+QsZOmzjilWux4CRMjer1mrRB04BmMwE9RCSkZOR9bVRAMB00iHVd9DyNOABYAoMtQvhSI+w0nbVIFNwNPbwA2M2ELENk+YwBuPn8aqzkc5Ujo-K1kYyK9d0MvLgB2KozauW4mluC27Nyo9UdesmLE4e8AVgnWrJnmyYXw5T0YKABzeCJQADMAJwhxJDIQWSQdEAALGDooNiQMCsOwgR7PJD9d4QJA+ai-f6AgggyLgl6IFLUD6IL4MLDAthQOhwX4A0FopCY6Gw6h4gleIkki6de5PdHbLEwxBvOllRmkllgtlIDnUxBwkC8wnEgXk4WIcqc160-F8mXMuUQhVKxBQqUM9Vk1HyxViqkIgFeIEgihKIA
\begin{tikzcd}
    {w_1: [x \mapsto 1]} \arrow[r, no head]                                                  & {w_2: [x\mapsto2]}                                                   &                                                                         \\
    {w_5: [x \mapsto 1; y \mapsto 3]} \arrow[u, dashed] \arrow[d, dashed] \arrow[r, no head] & {w_6: [x \mapsto 2; y\mapsto 4]} \arrow[u, dashed] \arrow[d, dashed] & {w_7: [x \mapsto 1; y\mapsto 4]} \arrow[llu, dashed] \arrow[ld, dashed] \\
    {w_3: [y \mapsto 3]} \arrow[r, no head]                                                  & {w_4: [y \mapsto 4]}                                                 &                                                                        
    \end{tikzcd}
    
New definition of the $\boxtimes$ modality:
\begin{itemize}
    \item Minimal world: a world $w$ is minimal for a property $\phi$ iff $\Dom{w} = \Code{FV}(\phi)$.
    \item $\boxtimes$ modality: $w \models \boxtimes \phi$ iff for all minimal worlds $w'$ that satisfy $\phi$, there exists a world $w''$ such that $w \sim w''$ and $w' \subseteq w''$.
\end{itemize}

\end{frame}

\begin{frame}[fragile]
    \frametitle{Better definition}

% https://tikzcd.yichuanshen.de/#N4Igdg9gJgpgziAXAbVABwnAlgFyxMJZABgBpiBdUkANwEMAbAVxiRAHcB9ARkQAJkADz4AdEQFs6aODgh9uFEAF9S6TLnyEU3clVqMWbLgCZ+QsZOmzjilWux4CRMjer1mrRB04BmMwE9RCSkZOR9bVRAMB00iHVd9DyNOABYAoMtQvhSI+w0nbVIFNwNPbwA2M2ELENk+YwBuPn8aqzkc5Ujo-K1kYyK9d0MvLgB2KozauW4mluC27Nyo9UdesmLE4e8AVgnWrJnmyYXw5T0YKABzeCJQADMAJwhxJDIQWSQdEAALGDooNiQMCsOwgR7PJD9d4QJA+ai-f6AgggyLgl6IFLUD6IL4MLDAthQOhwX4A0FopCY6Gw6h4gleIkki6de5PdHbLEwxBvOllRmkllgtlIDnUxBwkC8wnEgXk4WIcqc160-F8mXMuUQhVKxBQqUM9Vk1HyxViqkIgFeIEgihKIA
\begin{tikzcd}
    {w_1: [x \mapsto 1]} \arrow[r, no head]                                                  & {w_2: [x\mapsto2]}                                                   &                                                                         \\
    {w_5: [x \mapsto 1; y \mapsto 3]} \arrow[u, dashed] \arrow[d, dashed] \arrow[r, no head] & {w_6: [x \mapsto 2; y\mapsto 4]} \arrow[u, dashed] \arrow[d, dashed] & {w_7: [x \mapsto 1; y\mapsto 4]} \arrow[llu, dashed] \arrow[ld, dashed] \\
    {w_3: [y \mapsto 3]} \arrow[r, no head]                                                  & {w_4: [y \mapsto 4]}                                                 &                                                                        
    \end{tikzcd}
    
Examples:
\begin{itemize}
    \item $w_1 \models \boxtimes (1 \leq x \leq 2)$ holds.
    \item $w_5 \models \boxtimes (1 \leq x \leq 2)$ also holds. More generally, if $w \subseteq w'$ and $w \models \phi$, then $w' \models \phi$.
    \item $w_6 \models \boxtimes (1 \leq x \leq 2)$ also holds. More generally, if $w \sim w'$ and $w \models \boxtimes\phi$, then $w' \models \boxtimes\phi$.
\end{itemize}

\end{frame}

\begin{frame}[fragile]
    \frametitle{Separation}

% https://tikzcd.yichuanshen.de/#N4Igdg9gJgpgziAXAbVABwnAlgFyxMJZABgBpiBdUkANwEMAbAVxiRAHcB9ARkQAJkADz4AdEQFs6aODgh9uFEAF9S6TLnyEU3clVqMWbLgCZ+QsZOmzjilWux4CRMjer1mrRB04BmMwE9RCSkZOR9bVRAMB00iHVd9DyNOABYAoMtQvhSI+w0nbVIFNwNPbwA2M2ELENk+YwBuPn8aqzkc5Ujo-K1kYyK9d0MvLgB2KozauW4mluC27Nyo9UdesmLE4e8AVgnWrJnmyYXw5T0YKABzeCJQADMAJwhxJDIQWSQdEAALGDooNiQMCsOwgR7PJD9d4QJA+ai-f6AgggyLgl6IFLUD6IL4MLDAthQOhwX4A0FopCY6Gw6h4gleIkki6de5PdHbLEwxBvOllRmkllgtlIDnUxBwkC8wnEgXk4WIcqc160-F8mXMuUQhVKxBQqUM9Vk1HyxViqkIgFeIEgihKIA
\begin{tikzcd}
    {w_1: [x \mapsto 1]} \arrow[r, no head]                                                  & {w_2: [x\mapsto2]}                                                   &                                                                         \\
    {w_5: [x \mapsto 1; y \mapsto 3]} \arrow[u, dashed] \arrow[d, dashed] \arrow[r, no head] & {w_6: [x \mapsto 2; y\mapsto 4]} \arrow[u, dashed] \arrow[d, dashed] & {w_7: [x \mapsto 1; y\mapsto 4]} \arrow[llu, dashed] \arrow[ld, dashed] \\
    {w_3: [y \mapsto 3]} \arrow[r, no head]                                                  & {w_4: [y \mapsto 4]}                                                 &                                                                        
    \end{tikzcd}

Now we compose two equivalence classes $\{w_1, w_2\}$ and $\{w_3, w_4\}$ in a complete way:
\begin{itemize}
    \item $W_1$ and $W_2$ are disjoint equivalence classes, that is,
    \begin{align*}
        \forall w_1 \in W_1, \forall w_2 \in W_2, \Dom{w_1} \cap \Dom{w_2} = \emptyset
    \end{align*}
    \item $W_1 \oplus W_2$ is defined as the set of worlds:
    \begin{align*}
        \{ w_1 \cup w_2 ~|~ w_1 \in W_1 \land w_2 \in W_2 \}
    \end{align*}
    and all worlds in $W_1\oplus W_2$ are equivalent to each other. We also use $w_1 \oplus w_2$ to indicate a world $w_1 \cup w_2$ as defined above.
\end{itemize}

\end{frame}


\begin{frame}[fragile]
    \frametitle{Separation}

Consider a model $M_2$.

% https://tikzcd.yichuanshen.de/#N4Igdg9gJgpgziAXAbVABwnAlgFyxMJZABgBpiBdUkANwEMAbAVxiRAHcB9ARkQAJkADz4AdEQFs6aODgh9uFEAF9S6TLnyEU3clVqMWbLgCZ+QsZOmzjilWux4CRMjer1mrRB04BmMwE9RCSkZOR9bVRAMB00iHVd9DyNOABYAoMtQvhSI+w0nbVIFNwNPbwA2M2ELENk+YwBuPn8aqzkc5Ujo-K1kYyK9d0MvLgB2KozauW4mluC27Nyo9UdesmLE4e8AVgnWrJnmyYXwzrzVoh8BkqSRzgAOZT0YKABzeCJQADMAJwhxJBkECyJA6EAACxgdCgbEgYFYdhAv3+SH6wIgSCuEKhMK8cIRkWRAMQKWoIMQYIYWHhbCgdDgkJhiKJSFJ6Mx1CpNK8dIZLzOSL+xO2ZIxiCBXLKvMZApZiBF7MQWMltPpMuZQqQ5VFgM51Klav5GpRiG1irRKp5hqZhM1pp1JOokOhsIIBO+drZ5IVztx4DdsrtCvJoydONdNKUFCUQA
\begin{tikzcd}
    {w_1: [x \mapsto 1]} \arrow[r, no head]                                                  & {w_2: [x\mapsto2]}                                                                      &                                                                                            &     \\
    {w_5: [x \mapsto 1; y \mapsto 3]} \arrow[u, dashed] \arrow[d, dashed] \arrow[r, no head] & {w_6: [x \mapsto 2; y\mapsto 4]} \arrow[u, dashed] \arrow[d, dashed] \arrow[r, no head] & {w_7: [1; 4]} \arrow[llu, dashed] \arrow[ld, dashed] \arrow[r, no head] & {w_8: [2; 3]} \\
    {w_3: [y \mapsto 3]} \arrow[r, no head]                                                  & {w_4: [y \mapsto 4]}                                                                    &                                                                                            &    
    \end{tikzcd}

Examples:
\begin{itemize}
    \item $w_1 \oplus w_3$ (i.e., $w_5$) $\models \boxtimes (1 \leq x \leq 2) * \boxtimes (3 \leq y \leq 4)$
    \item $w_1 \oplus w_3$ (i.e., $w_5$) $\models (x : \urt{\Int}{1 \leq \vnu 2}) * (y : \urt{\Int}{3 \leq \vnu 4})$
    \item $w_4 \models \boxtimes (x = 2) \wand \boxtimes (y = x + 2)$
    \item Also, $w_3 \models \boxtimes (x = 2) \wand \boxtimes (y = x + 2)$
    \item For all equivalence classes $W$ that make $\boxtimes(x = 2)$ hold and $W \cap \{w_3, w_4\} = \emptyset$, we have $W \oplus \{w_3, w_4\}$ makes $\boxtimes (y = x + 2)$ hold (since $w_3 \sim w_4$, and $x = 2 \implies y = x + 2$ in $w_4$). 
\end{itemize}

\end{frame}

\begin{frame}[fragile]
    \frametitle{Modeling Context}

We use Kripke semantics to model context; that is, a (type) context is a set of models, where each model contains multiple equivalence classes.

\textbf{Model:} A model is a tuple $(W, \sim, \sqsubseteq, w)$, where $W$ is a set of worlds, $\sim$ is an equivalence relation over $W$, $\sqsubseteq$ is a partial ordering over $W$, and $w$ is a world in $W$.

\begin{itemize}
    \item $x{:}\ort{\Int}{\vnu > 0}$ can be interpreted as $x > 0$, where
    \begin{align*}
        \denot{x > 0} \doteq \{ (\{[x\mapsto v]\}, \emptyset, \emptyset, [x\mapsto v]) ~|~ v > 0 \}
    \end{align*}
    \item Each model contains a single world; there is neither an equivalence nor a partial ordering relation.
\end{itemize}


\end{frame}

\begin{frame}[fragile]
    \frametitle{Modeling Context}

\textbf{Model:} A model is a tuple $(W, \sim, \sqsubseteq, w)$, where $W$ is a set of worlds, $\sim$ is an equivalence relation over $W$, $\sqsubseteq$ is a partial ordering over $W$, and $w$ is a world in $W$.

\begin{itemize}
    \item $x{:}\urt{\Int}{\vnu > 0}$ can be interpreted as $\boxtimes (x > 0)$, where
    \begin{align*}
        \denot{\boxtimes (x > 0)} \doteq \{ (\{[x\mapsto v] ~|~ v > 0\}, \_, \emptyset, [x\mapsto 1]) \}
    \end{align*}
    \item There is only one model, which contains all possible worlds $[x\mapsto v]$ where $v$ is a positive integer; all of these worlds are equivalent to each other.
\end{itemize}


\end{frame}

\begin{frame}[fragile]
    \frametitle{Modeling Context}

\textbf{Model:} A model is a tuple $(W, \sim, \sqsubseteq, w)$, where $W$ is a set of worlds, $\sim$ is an equivalence relation over $W$, $\sqsubseteq$ is a partial ordering over $W$, and $w$ is a world in $W$.

\begin{itemize}
    \item $x{:}\urt{\Int}{1 \leq \vnu \leq 2} * y{:}\urt{\Int}{3 \leq \vnu \leq 4}$ can be interpreted as $\boxtimes (1 \leq x \leq 2) * \boxtimes (3 \leq y \leq 4)$.
    \item There is a single model (i.e., $M_2$), and all worlds are equivalent.
    \item $x{:}\urt{\Int}{1 \leq \vnu \leq 2},\; y{:}\urt{\Int}{3 \leq \vnu \leq 4}$ can be interpreted as $\boxtimes (1 \leq x \leq 2) \land \boxtimes (3 \leq y \leq 4)$.
    \item There are multiple models to consider; $M_1$ and $M_2$ are examples of them.
\end{itemize}

\end{frame}

\begin{frame}[fragile]
    \frametitle{Operational Semantics under Context}

Remember that we use Kripke semantics to express aligned multiple executions together. If we don't align them, we can simply express them as $\mathcal{P}(\mathcal{P}(\sigma))$.

\textbf{Intuition:} A model can be extracted back to a set of substitutions (equivalence class).

\textbf{Example:}
\begin{align*}
&\eraserf{M_1} = \{ [x\mapsto 1; y\mapsto 3], [x\mapsto 2; y\mapsto 4] \}
\\&\eraserf{M_2} = \{ [x\mapsto 1; y\mapsto 3], [x\mapsto 1; y\mapsto 4], [x\mapsto 2; y\mapsto 3], [x\mapsto 2; y\mapsto 4] \}
\end{align*}

\textbf{Reduction under a model:}
\begin{align*}
    (W, \sim, \subseteq, w) \models \mstep{e}{v} \iff \exists w' \in W.\; w \sim w' \land \mstep{w'(e)}{v}
\end{align*}

\end{frame}

\begin{frame}[fragile]
    \frametitle{Type denotation}

\textbf{Base Type:}
\begin{align*}
    \denot{\ort{b}{\phi}}_M &\doteq \{ e ~|~ \forall v. M \models \mstep{e}{v} \implies \phi[\vnu \mapsto v] \} \\
    \denot{\urt{b}{\phi}}_M &\doteq \{ e ~|~ \forall v. \phi[\vnu \mapsto v] \implies M \models \mstep{e}{v} \}
\end{align*}

\textbf{Function Type:}
\begin{align*}
    \denot{x{:}\tau_x \sarr \tau}_M &\doteq \{ e ~|~ \forall v_x \in \denot{\tau_x}_M. (e\;x) \in \denot{\tau}_{M[x\mapsto v_x]} \} \\
    \denot{x{:}\tau_x \wand \tau}_M &\doteq \{ e ~|~ \forall M'.\; \denot{x : \tau_x}_M.\; \Dom{M'} \cap (\Dom{M} \cup \Code{FV}(\tau)) = \emptyset \implies 
    \\&\qquad (e\;x) \in \denot{\tau}_{M \oplus M'} \}
\end{align*}
where $M[x\mapsto v_x]$ means to add $[x\mapsto v_x]$ to all worlds in $M$.

\vspace{0.5cm}
\textbf{Fundamental Theorem:}
\begin{align*}
    \Gamma \vdash e : \tau \implies \forall M.\; M \in \denot{\Gamma}.\; e \in \denot{\tau}_M
\end{align*}

\end{frame}

\begin{frame}[fragile]
    \frametitle{Type denotation}

\textbf{Function Type Examples:}

\begin{align*}
    &\denot{x{:}\urt{b}{1 \leq \vnu \leq 2} \wand \urt{b}{3 \leq \vnu \leq 4}}_{\emptyset} 
    \\\doteq &\{ e ~|~ \forall M'.\; \denot{\boxtimes (1 \leq x \leq 2)}_M.\; \Dom{M'} \cap (\Dom{M} \cup \Code{FV}(\tau)) = \emptyset \implies 
    \\& (e\;x) \in \denot{\tau}_{M \oplus M'} \}
    \\&\text{Consider one model}~ M_0 = (\{[x\mapsto 1], [x\mapsto 2]\}, \sim, \emptyset, [x\mapsto 1])
    \\\doteq &\{ e ~|~ (e\;x) \in \denot{\urt{b}{3 \leq \vnu \leq 4}}_{M_0} \}
    \\ \doteq &\{ e ~|~ M_0 \models \mstep{e\;x}{3} \land M_0 \models \mstep{e\;x}{4} \}
    \\ \doteq &\{ e ~|~ (\mstep{e\;1}{3} \lor \mstep{e\;2}{3}) \land (\mstep{e\;1}{4} \lor \mstep{e\;2}{4}) \}
\end{align*}
Which is consistent with incorrectness triples.

\end{frame}